
déjà
prouvé, on en conclut que quelque soit \(m\ge -2\),
\[
h^1(L(n))\le
\bigl[2+h_2(a_{r-2},a_{r-1},a_r)+m\bigr]
h_{r-1}(a_{r-2},a_{r-1},a_r)
\tag{2.9.2}
\]
pour tout \(n\le m\).

Par 2.5, \(L_Y\) est un \((0,\ldots,0,a_r)\)-faisceau. Soient \(c_i=g_{r-i}(a_{r-2},a_{r-1},a_r)\) pour \(i=1,\ldots,r\) et \(m=P_{r-2}(c_1,\ldots,c_{r-1})\). Par 1.11, \(m\ge0\) et \(L_Y\) est \(m\)-régulier. Donc par 1.4 (ii), \(h^1(L(m-2))\ge h^1(L(n))\) pour tout \(n\ge m-1\). Donc si l'on définit des polynômes par \(M=P_{r-2}(g_{r-1},\ldots,g_1)\) et \(h_r=(h_2+M)h_{r-1}\), alors on a (ii) en vertu de 2.9.2.

Enfin, comme \(L_Y\) est \(m\)-régulier, \(h^q(L(m))=0\) pour \(q\ge2\) par 1.4 (i)\,; d'où \(\sum_{i=0}^{r}a_i\binom{m+i}{i}=h^0(L(m))-h^1(L(m))\). Mais \(h^0(L(m))\le a_r\binom{m+r}{r}\) par 1.8 (i). Donc si l'on définit \(g_r=\sum_{i=1}^{r}g_{r-i}\binom{M+i}{i}+a_r\binom{M+r}{r}+h_r\), on a (iii).

\textbf{Remarque 2.10.} Pour \(m\) fixé, considérons une égalité de polynômes en \(n\) :
\[
\sum_{i=0}^{r}a_{m,i}\binom{n+i}{i}
=
\sum_{j=0}^{r}a_j\binom{m+n+j}{j}.
\]
Alors on trouve facilement la formule
\[
a_{m,i}=\sum_{j=0}^{r-i}a_{j+i}\binom{m-1+i}{j}
\]
en raisonnant par récurrence sur \(r\) afin que \(a_{m,i}\) devienne \(a_{m,0}\), et en faisant \(n=-1\).

De même, on trouve qu'il existe un polynôme à coefficients rationnels \(P_i^{(m,r)}(X_r,\ldots,X_i)\) qui est croissant en \(X_i\), tel que dans l'égalité suivante :
\[
\sum_{i=0}^{r}a_i^{(m)}\binom{n+i}{i}
=
\sum_{j=0}^{r}a_j\binom{mn+j}{j}
\]
